\section{Discussion}

Our results demonstrate that Pilot-Driven Viability Gates achieve 93.3\% accuracy for early identification of computationally infeasible hypotheses using synthetic validation. We interpret these findings, acknowledge limitations transparently, and assess broader implications for ML research methodology.

\subsection{Key Findings}

\textbf{Finding 1: Perfect linear scaling is a synthetic artifact, but the mechanism is sound.}

The r=1.000 correlation between micro-pilot and full-scale overhead validates the framework's extrapolation logic in idealized conditions. However, this perfect linearity is unlikely in real experiments. We expect three sources of deviation:

\begin{enumerate}
\item \textbf{Measurement noise}: Hardware variance, system load, I/O contention introduce $\pm$5--10\% timing fluctuations. Real correlation expected r=0.7--0.9.

\item \textbf{Non-linear scaling}: Memory bottlenecks (100-sample dataset fits in cache, full dataset spills to RAM) or I/O overhead (data loading dominates for large datasets) create non-linear scaling patterns not captured by $O_{\text{full}} = k \times O_{10}$.

\item \textbf{Type-specific scaling}: Attention mechanisms (O(n$^2$) in sequence length) may exhibit different $k$ than gradient penalties (O(n)). Our CV=0.00\% across types is artifact; real data expected CV=10--30\%.
\end{enumerate}

Despite these expected deviations, the $r \geq 0.7$ threshold provides buffer. Even with measurement noise reducing correlation from r=1.000 to r=0.75, the mechanism remains valid. The risk is $r < 0.7$, where extrapolation breaks down---this failure mode should trigger framework recalibration (e.g., memory profiling for non-linear cases).

\textbf{Finding 2: Marginal Bayesian benefit reflects strong prior, not mechanism failure.}

Error reduction 40.91\% vs 40\% threshold (0.91pp excess) might suggest fragile benefit. However, the paired t-test (p=0.0003) confirms statistical significance, and absolute error reduction (84\%: 0.6966 $\rightarrow$ 0.1111) is substantial. The marginal relative reduction reflects synthetic data's strong prior (k=1.000 perfect scaling leaves little room for Bayesian refinement).

Real-world scenarios with weaker correlation (r=0.7--0.8) will have higher Gate 1 prior error, providing more opportunity for Gate 2 likelihood to refine predictions. We hypothesize error reduction >50\% in real data. The mechanism works as designed; the marginal result is artifact of idealized conditions.

\textbf{Finding 3: High recall (96.2\%) on non-viable hypotheses minimizes false negatives.}

The framework's design priority---minimize missed non-viable hypotheses---is empirically validated. Only 1 of 26 non-viable hypotheses escaped Gate 1 (false negative rate 3.8\%). This false negative occurred at 10.3\% overhead, barely above the 10\% threshold, suggesting measurement uncertainty rather than systematic prediction failure.

The single false positive (viable hypothesis incorrectly stopped) occurred at 9.7\% overhead, also borderline. For non-borderline cases (overhead <9\% or >11\%), classification was 100\% accurate. This suggests incorporating confidence intervals: borderline cases (within $\pm$1\% of threshold) should proceed to Gate 2 for refined prediction rather than stopping at Gate 1.

\subsection{Limitations}

We document limitations transparently to guide interpretation and future work. These constraints bound the scope of valid claims.

\textbf{L1: Synthetic Corpus Limits External Validity}

Our validation used synthetic data with perfect linear scaling (k=1.000, r=1.000, CV=0.00\%). This establishes proof-of-concept---the framework's gate logic, statistical tests, and Bayesian updates work as designed. However, external validity is unknown: Does the $r \geq 0.7$ correlation hold on real published papers? Do real hypothesis types exhibit consistent scaling (CV < 30\%)?

\textbf{Boundary Condition}: Framework validated in synthetic contexts with perfect linear scaling. Real-world applicability contingent on real corpus validation showing $r \geq 0.7$ (High Priority Future Work FD1).

\textbf{L2: Single Threshold Tested}

All validation used T=10\% overhead threshold (real-time deployment constraint). Framework behavior at permissive thresholds (T=50\% for offline batch processing, T=200\% for research contexts) is unknown. The 87\% non-viable prevalence at 10\% threshold may shift to 50\% at 50\% threshold or 10\% at 200\% threshold, changing accuracy characteristics.

\textbf{Boundary Condition}: Framework validated for strict deployment thresholds (10\%). Multi-threshold validation needed to establish threshold-accuracy curve (Medium Priority Future Work FD2).

\textbf{L3: Type-Specific Scaling Untested}

Synthetic corpus enforced global k=1.000 across all hypothesis types (attention, gradient, regularization, normalization). Real data may require per-type calibration if scaling variance exceeds CV=30\% threshold. This would increase framework complexity (maintain $k_{\text{attention}}$, $k_{\text{gradient}}$ lookup table) but improve accuracy.

\textbf{Boundary Condition}: Framework validated assuming $k$ generalizes (Assumption A3). If real data shows CV > 30\%, per-type calibration required (Low Priority Future Work FD5).

\textbf{L4: User Compliance Untested}

Gate 1 achieved 93.3\% \textit{classification accuracy} (can identify non-viable hypotheses). Assumption A4 posits researchers will \textit{act on stop signals} (compliance behavior). This is untested. Potential failure modes:

\begin{itemize}
\item Confirmation bias: Researcher ignores Gate 1 stop signal for favored hypothesis
\item Sunk cost fallacy: Prior investment in hypothesis design motivates full implementation despite stop signal
\item Publication pressure: Negative results (stopped hypotheses) harder to publish than positive results
\end{itemize}

\textbf{Boundary Condition}: Framework validated as \textit{predictive tool} (93.3\% accuracy). Adoption as \textit{decision-making tool} requires prospective user study measuring stop rate $\geq$60\%, time savings $\geq$3 hours (High Priority Future Work FD3).

\textbf{L5: Imbalanced Corpus May Inflate Accuracy}

With 26 non-viable and 4 viable hypotheses (87\% non-viable prevalence), the 93.3\% accuracy metric is dominated by high recall (96.2\%) on the large non-viable class. The false positive rate (25\%: 1/4 viable stopped) is less representative due to small viable sample size.

Balanced corpus validation (15 viable, 15 non-viable at 50\% threshold) would test whether accuracy $\geq$80\% holds with equal class distribution. We anticipate accuracy reduction to 85--90\% under balanced prevalence but still above 80\% threshold.

\textbf{Boundary Condition}: Framework validated under imbalanced distribution (87\% non-viable). Balanced validation recommended to confirm accuracy $\geq$80\% with 50--50 prevalence.

\subsection{Broader Impact}

\textbf{Positive Impacts}:

\begin{enumerate}
\item \textbf{Reduced research waste}: h-e1 anecdote shows 68.65\% overhead discovered after full implementation. Framework enables micro-pilot detection (<1 hour) before days of wasted effort.

\item \textbf{Systematic methodology}: Fills gap in ML research workflow. Ablation studies test variations; complexity analysis provides theoretical bounds; viability gates add empirical early-stop protocol.

\item \textbf{Quantified uncertainty}: Bayesian posterior provides confidence intervals, not binary decisions. Researchers can assess borderline cases ($\pm$1\% of threshold) with probabilistic reasoning.
\end{enumerate}

\textbf{Negative Impacts and Mitigation}:

\begin{enumerate}
\item \textbf{False positives risk missed opportunities}: 1 of 4 viable hypotheses incorrectly stopped (25\% FP rate on small sample). A viable hypothesis rejected at Gate 1 is a lost contribution.

   \textbf{Mitigation}: Borderline cases (confidence interval straddles threshold) proceed to Gate 2 for refined prediction. Only high-confidence non-viable predictions ($O_{\text{pred}} > 1.1T$) trigger immediate stop.

\item \textbf{Overreliance on micro-pilot}: Researchers may skip hypothesis refinement, trusting Gate 1 predictions uncritically.

   \textbf{Mitigation}: Framework provides probability distributions and confidence intervals, encouraging critical evaluation rather than blind acceptance.

\item \textbf{Publication bias amplification}: Stopped hypotheses yield no positive results, reducing negative result publication.

   \textbf{Mitigation}: Framework documentation (honest reporting protocol) encourages reporting Gate 1 stops as methodological contribution, not failure.
\end{enumerate}

\textbf{Ethical Considerations}:

The framework assumes computational efficiency as a valued constraint. In resource-constrained settings (limited compute budgets, energy efficiency mandates), this assumption aligns with sustainability goals. However, premature stopping may hinder exploratory research where ``inefficient'' hypotheses later inspire efficient variants. We recommend balancing viability gates (feasibility-first) with exploratory investigation (curiosity-driven) rather than replacing the latter.

\subsection{Threats to Validity}

\textbf{Internal Validity}: Synthetic corpus provides controlled comparison (same hypotheses across all gates). No confounds from hardware variance, dataset differences, or implementation variations. Risk: Synthetic data may mask real-world complexities (memory bottlenecks, I/O overhead, non-linear scaling).

\textbf{External Validity}: The limiting threat. Synthetic validation establishes proof-of-concept but not generalizability. Real corpus validation (FD1) is critical next step. Until $r \geq 0.7$ confirmed on real published papers, claims scope to ``framework mechanics validated in synthetic conditions.''

\textbf{Construct Validity}: Accuracy, correlation, and error reduction measured correctly with standard statistical tests (binomial, Pearson, paired t-test). Time investment (``<1 hour'' for Gate 1) not empirically validated---synthetic runtime <1 second not representative of real micro-pilot cost. Prospective validation (FD3) needed to confirm time savings claim.

\textbf{Statistical Conclusion Validity}: All results show strong statistical significance ($p < 0.05$). Effect sizes large for M1 (r=1.000) and M3 (93.3\% accuracy), marginal for M2 (40.91\% vs 40\%). Sample sizes adequate: 32 hypotheses (M1), 20 hypotheses (M2), 30 hypotheses (M3). Risk: Perfect correlation (r=1.000) and zero variance (CV=0.00\%) unlikely to replicate in real data.

This transparent limitations analysis guides appropriate interpretation: Framework mechanics validated under synthetic conditions. External validity, multi-threshold behavior, user compliance, and balanced-corpus accuracy remain open questions for real-world validation.
